
\chapter{Conclusion}
\label{ch:conclusion}

\section{Deep Learning and LHC Physics}
\label{sec:conclusion-dl-lhc}

Machine learning is nothing new to LHC physics, but the deep learning revolution kicked off by AlexNet has produced a large suite of new tools that have transformed how LHC data is processed into physics results.
The research covered in this thesis has been a small part of this transformation.
\Cref{ch:tagging} presented a study of top-quark jet tagging using low-level constituent inputs, which established that neural networks trained directly on jet constituent kinematics outperform those trained on hand-engineered expert features, mirroring the lesson of AlexNet in the computer vision field.
The study also provided a systematic uncertainty quantification framework which showed that larger, more expressive taggers extract more information from simulation, including information that is imperfectly modeled, so raw performance gains in simulation do not trivially translate into performance gains on data.
This effect is visible in the most recent ATLAS and CMS flavor tagging results~\cite{ATLAS:2025dkv,Sarkar:2024vjz}, and managing it will be a challenge as jet tagging models continue to scale.

\Cref{ch:unfolding} surveyed the new generation of unbinned unfolding methods, including diffusion- and flow-matching-based generative approaches and the \Omnifold algorithm, which recasts unfolding as an iterative likelihood ratio estimation problem.
\Cref{ch:zjets} then demonstrated \Omnifold in practice, producing the highest-dimensional cross section measurement ever performed at a hadron collider.
The result is a reweighted particle-level dataset encoding the full differential cross section of $Z(\to\mu\mu)$+jets production in up to 843 simultaneous dimensions.
Unlike a binned measurement, whose output is a fixed set of histograms tied to the observables chosen at analysis time, this result is a dataset that can be projected onto any observable of interest, including observables that could not have been unfolded by any other means.
Four physics use cases of this dataset were presented, three of which would be very difficult or impossible to perform with traditional binned unfolding techniques, proving that deep-learning-based unfolding methods enable entirely new kinds of measurements.
I believe these new capabilities will produce a shift in how the community understands the purpose of cross section measurements.
Instead of targeting specific observables one measurement at a time, the focus will shift to making general-purpose measurements of various final states and prioritizing downstream applications.
This will allow for flexible use of the data, possibly many years into the future after the LHC is done running and the experimental collaborations no longer exist.

Jet tagging and unfolding were the deep learning applications covered in this thesis, but there are many more examples in LHC physics.
Anomaly detection, simulation-based inference, fast detector simulation, and event reconstruction have all seen substantial activity.
There are few parts of LHC physics that deep learning does not touch in some way.
Despite this progress, no major experimental discovery, for example new particles or dynamics beyond the Standard Model, has been produced.
LHC physics is in many ways still organized around the same measurement and search strategies it used before the deep learning era.
Individual pieces of the data analysis chain have been replaced or substantially improved by neural networks, but the overall structure of how physics results are produced has not yet changed.
I believe that the unbinned and high-dimensional measurements described in \Cref{ch:unfolding,ch:zjets} represent genuine progress in how particle physics measurements are performed, and that the order-of-magnitude improvements in jet classification documented in \Cref{ch:tagging} are producing real gains in the sensitivity of searches for rare processes like di-Higgs production, but all of this has not been enough to deliver progress on the open problems in \Cref{sec:open_problems}.
The sections that follow discuss the two directions I find most promising for changing this.

\section{Higgs and di-Higgs with Large-scale Flavor Tagging}
\label{sec:conclusion-dihiggs}

The Higgs boson is rightfully the current focus of the LHC physics program.
Given that $H \to b\bar{b}$ accounts for roughly 58\% of Higgs decays, flavor tagging is a very important ingredient of many Higgs studies.
This is especially true in searches for di-Higgs production and efforts to constrain the Higgs self-coupling parameter $\lambda$.
This parameter is not a free parameter of the Standard Model.
The Standard Model provides a prediction for $\lambda$ in terms of the other free parameters, and this prediction can be tested by measuring the di-Higgs production cross section.
The self-coupling $\lambda$ determines the shape of the Higgs potential, and any departure from the Standard Model prediction would shift that shape in ways that have significant consequences for many of the open problems listed in \Cref{sec:open_problems}, including the hierarchy problem and the matter-antimatter asymmetry.

The most sensitive di-Higgs production channels ($HH \rightarrow b\bar{b}\gamma\gamma$, $HH \rightarrow b\bar{b}\tau\tau$, etc.) all involve one of the Higgs bosons decaying to a bottom quark pair to increase the cross section.
As discussed in \Cref{sec:tagging-future}, the order-of-magnitude gains in $b$-jet rejection achieved by the current generation of transformer-based flavor taggers~\cite{ATLAS:2025dkv,Sarkar:2024vjz}, and the further gains projected from scaling these models to larger architectures trained on billions of jets~\cite{ATLAS:2026vyw,Vigl:2026ppx}, make an HL-LHC observation of di-Higgs production and a constraint on the Higgs self-coupling realistic.
This would be the flagship physics result of the HL-LHC, and continued work on scaling jet classification methods strikes me as essential to making it a reality.

Such a di-Higgs observation could go one of two ways.
If the measured self-coupling is consistent with the Standard Model prediction, it will be an extraordinary confirmation of a theory we already know to be incomplete, and the open problems listed in \Cref{sec:open_problems} will remain exactly as open as they are today.
The direction forward for particle physics in that case is unclear.
We would have a theory that fits all data with no clear experimental goal that could point the way forward.
If the self-coupling deviates from the Standard Model prediction, we will have the first direct evidence that the Higgs sector is non-minimal, pointing toward explanations of the hierarchy problem, the matter-antimatter asymmetry, or both.
In either case, a di-Higgs observation will be an enormous achievement.
Deep learning, including the flavor tagging results of \Cref{ch:tagging}, will be one of the many innovations that make it possible.

\section{Agentic AI and Physics Research}
\label{sec:conclusion-agentic}

Almost all research on deep learning applications in particle physics has used it as a tool for processing data.
Recently the emergence of agentic AI systems has made it possible to use deep learning to perform the research process itself, rather than simply process data.
Agentic systems can plan multi-step analyses, execute code, search literature, and iterate on their own outputs.
As of May of 2026, such systems can already reproduce existing particle physics analyses with high fidelity~\cite{Birk:2026zpd} and produce reasonable, if imperfect, results on open-ended data analysis tasks~\cite{Moreno:2026mqk}.
If model capability stopped improving today, I think this new capability would already be transformative.

Within the context of particle physics data analysis, the single most constrained resource in experimental collaborations is human time.
The cross section measurement of \Cref{ch:zjets} is only possible because of the dedicated effort of hundreds of researchers who operated the ATLAS detector during run 2, designed analysis software, calibrated analysis objects, and set systematic uncertainties on track measurements.
Maintaining and improving the complicated pipeline that turns raw collision data into physics results is unglamorous but extremely important work.
Typically, the methodology of such work is well understood, but software implementation can be difficult and frustrating.
Furthermore, many of the working groups responsible for these tasks are critically understaffed in ATLAS, making the downstream physics analyses suboptimal.
Agentic AI has the potential to automate these tasks and eliminate the human time bottleneck.
Particle physics experiments could become dramatically more productive if agentic AI can follow up on all possible improvements and deliver a level of optimality that is impossible to support with the current human time constraints.

Additionally, I believe the LHC datasets contain far more physics than human researchers have been able to extract from them.
To give a few examples, there are large regions of unexplored two-body resonance parameter space~\cite{Craig:2016rqv,Kim:2019rhy}, there are still gaps in parameter space for many popular BSM models, and many interesting deep-learning-based data analysis techniques have never been utilized.
The LHC datasets support all of this physics, but there are insufficient human researchers willing to write the necessary software, especially in cases where a research direction is only partially complete but the community has by and large lost interest.
Agentic AI can fill these gaps, without the time and career constraints that limit human scientists.

The current paradigm of fundamental physics research features a division of labor between humans and AI systems, where the AI handles the computational ``backend'' of research while the human handles the ideation, high-level implementation, and communication of results.
This is often called ``centaur science'', in reference to the mythical creature that is upper half human and lower half horse.
In my opinion, this is a remarkably exciting moment to be a human researcher.
Agentic AI has liberated us from the need to spend the majority of our working hours developing and maintaining software, shifting the bottleneck from implementation to creativity.
I predict that much more elaborate and inventive data analyses will become standard in the near-term future.
To give a concrete example, given the computational complexity of the intrinsic dimension measurement in \Cref{sec:zjets-results-id}, I could not have produced those results in the timeframe of my Ph.D. without agentic coding tools.
I was responsible for the ideation and high-level methodology, but the necessary software development and optimization was entirely performed by agentic tools.

My expectation is that the current balance between human researcher and AI will not last long.
Particle physics data analysis tasks are, at their core, software engineering and logical inference problems, both areas in which AI agents have shown rapid improvement.
As of May of 2026, AI systems are still not as good as a competent human researcher at extracting physics from data, but the gap is not very wide and the systems will only get better in the future.
Over the past decade, betting against all but the most aggressive predictions about AI capability has been a losing proposition.
I think it is likely that within a few years we will have AI systems that can take a physics question, identify the relevant dataset, design and execute the analysis, interpret the result, and propose the next step.
Setting academic labor market implications aside, I believe such a system could produce extremely rapid progress.
A few thoughts on what this may look like, and the conclusion of the thesis, are offered in the next section.

\section{Will AI Deliver the Next Discovery?}
\label{sec:conclusion-discovery}

My prediction is that the combination of HL-LHC data and AI-powered analysis will allow us to essentially exhaust the statistical potential of all available experimental data within the next decade.
I think we are currently operating far below this level, with the potential to make much more accurate claims about the Standard Model and $\Lambda$~CDM cosmology currently or soon to be sitting on disk.
In particle physics specifically, the painstaking accumulation of results from the first three runs of the LHC has brought us closer to exhausting the potential of the data than any previous moment, and this thesis has been a small contribution to that effort.
However, I suspect the remaining gap is wide, and AI will close this gap to near zero in the coming years.

This will produce one of two outcomes.
The first possibility is that the Standard Model is confirmed to be the correct effective theory of the electroweak scale, including a measurement of the Higgs self-coupling consistent with the SM prediction.
This would be a remarkable scientific achievement, but would also be deeply unsatisfying since the Standard Model is known to be incomplete.
In this case, the field of particle physics would be forced to look toward collecting new and different types of experimental data in the hope that some of them would provide clues to answer the open questions of \Cref{sec:open_problems}.
Neutrino experiments, cosmological probes, cosmic ray observatories, or a next-generation collider would be obvious avenues.

The second possibility is that AI-powered analysis of the HL-LHC data surfaces a deviation from Standard Model predictions.
If this happens, we may produce real answers to some of the open questions in particle physics.
Such a discovery would hopefully make obvious what future experimental efforts to prioritize, but it is likely that the availability of experimental data will always remain a bottleneck.
I think it is unlikely that a superintelligent AI system will simply return the correct Lagrangian of the universe.

I am personally excited about both of these outcomes, and the prospect of knowing which one will occur in a matter of years rather than decades.
Hopefully the work described in this thesis will be part of the foundation on which these largely AI-driven data analyses will be built.
I am optimistic about the probability of the next decade producing tangible progress in particle physics, which by all accounts is badly needed.
Developing and deploying the tooling of modern LHC data analysis has been a fulfilling and fruitful research direction, and I am excited about the potential of AI to accelerate this research.
The day-to-day activity of doing particle physics research is currently changing faster than at any point in the history of the field, and this is typically a promising sign that change is on the horizon.
My hope is that all of this change will teach us something new.

\FloatBarrier
